\chapter{Онолын судалгаа ба ижил төстэй системүүд}

\section{Компьютерын харааны тухай ойлголт, үүсэл хөгжил}
Компьютерын хараа (Computer Vision) нь хиймэл оюун ухааны нэг гол салбар бөгөөд дижитал зураг эсвэл бичлэгээс утга учир бүхий мэдээллийг автоматаар гаргаж авах, боловсруулах, таних зорилготой шинжлэх ухаан юм. Хүн төрөлхтний нүд болон тархи дүрсэн мэдээллийг хэрхэн хүлээж авч ойлгодогтой ижил байдлаар компьютерыг "харах" чадвартай болгоход чиглэдэг. Эхэн үеийн компьютерын харааны аргууд нь дүрсийн ирмэг, өнгө, хэлбэр зэрэг гар аргаар тодорхойлсон шинж чанарууд (hand-crafted features)-д тулгуурладаг байв (жишээ нь: SIFT, HOG, SURF). Эдгээр нь бага хэмжээний өөрчлөлттэй орчинд сайн ажилладаг байсан боловч гэрэлтүүлэг өөрчлөгдөх, дүрс хагас халхлагдах үед таних нарийвчлал маш ихээр буурч байсан тул нарийн төвөгтэй бодит орчин дахь асуудлуудыг бүрэн шийдэж чаддаггүй байв. 

2010-аад оноос хойш Гүнгэрвэл сургалт (Deep Learning) болон Конволюцийн Мэдрэлийн Сүлжээ (Convolutional Neural Networks - CNN)-ийн хөгжил үсрэнгүй урагшилснаар компьютерын харааны салбарт жинхэнэ хувьсгал гарсан. CNN нь оролтын зургаас багаас их рүү чиглэсэн хэв шинжүүдийг (low-level to high-level features) олон давхаргуудын тусламжтайгаар автоматаар суралцдаг ба энэ нь өмнөх үеийн уламжлалт аргуудаас хамаагүй илүү нарийвчлалтай болж чадсан юм.

\section{Объект илрүүлэлтийн алгоритмуудын дэвшил}
Компьютерын харааны нэгэн чухал дэд салбар болох Объект илрүүлэлт (Object Detection) нь тухайн зурган дээр байгаа объектуудын зөвхөн төрлийг (classification) хэлээд зогсохгүй, тэдгээр объектууд зургийн хаана байрлаж байгааг тэгш өнцөгт хайрцгаар (bounding box) нарийвчлан зааж өгөх үйлдлийг багтаадаг. Объект илрүүлэлтийн алгоритмууд нь хөгжлийнхөө явцад үндсэн хоёр том бүлэгт хуваагдан хөгжиж ирсэн. Үүнд: 1) Хоёр алхмаар илрүүлдэг алгоритмууд (Two-stage detectors), 2) Нэг алхмаар илрүүлдэг алгоритмууд (One-stage detectors).

\subsection{Хоёр алхмаар илрүүлдэг алгоритмууд (Two-stage Detectors)}
Эдгээр загварууд нь илрүүлэлтийг хоёр тусдаа үе шаттайгаар гүйцэтгэдэг учраас нарийвчлал өндөртэй ч, ажиллах хурд удаан байдаг шинж чанартай.
\begin{itemize}
    \item \textbf{R-CNN (Region-based CNN):} 2014 онд танилцуулагдсан бөгөөд зурган дээр объект байж болзошгүй 2000 орчим бүсүүдийг (region proposals) урьдчилан ялгаж, бүс тус бүр дээр CNN ашиглан объектын онцлогийг сургаж таньдаг аргачлалтай. Энэ нь маш их хугацаа шаарддаг байв.
    \item \textbf{Fast R-CNN:} R-CNN-ийн хурдны асуудлыг шийдэж, бүх зургийг нэг л удаа CNN дээгүүр явуулж, түүнээс гаргаж авсан feature map дээрээс бүсүүдийг ялгах болсноор хурдассан боловч region proposal-ийг хөндлөнгийн Selective Search алгоритм ашиглаж хийдэг хэвээр байв.
    \item \textbf{Faster R-CNN:} Selective Search-ийг Region Proposal Network (RPN) хэмээх жижиг мэдрэлийн сүлжээгээр сольж өгснөөр бүрэн утгаараа end-to-end сургах боломжтой болсон алгоритм юм \cite{faster_rcnn}. Хэдийгээр хурд эрс сайжирсан ч бодит хугацаанд (real-time) өндөр давтамжтай видео дээр ажиллахад бэрхшээлтэй хэвээр байв.
\end{itemize}

\subsection{Нэг алхмаар илрүүлдэг алгоритмууд (One-stage Detectors)}
Эдгээр загварууд нь зургаас объект байрлах бүсийг түрүүлж хайлгүйгээр, зургийг шууд харж нэг л удаагийн үнэлгээгээр (single pass) объектын ангилал ба байршлыг нэгэн зэрэг таамагладаг онцлогтой. Энэ нь нарийвчлалыг бага зэрэг золиосолж, хурдыг асар ихээр нэмэгдүүлсэн технологи юм.
\begin{itemize}
    \item \textbf{SSD (Single Shot MultiBox Detector):} Төрөл бүрийн хэмжээтэй объектуудыг илрүүлэхийн тулд CNN-ийн өөр өөр давхаргуудын feature map-уудыг зэрэг ашигладгаараа онцлог.
    \item \textbf{YOLO (You Only Look Once):} 2015 онд Joseph Redmon болон бусад судлаачдын танилцуулсан YOLO алгоритм нь объект илрүүлэлтийн салбарт жинхэнэ хувьсгал хийсэн \cite{yolo_original}. YOLO нь оролтын зургийг $S \times S$ хэмжээтэй торон хүснэгтэд хуваах бөгөөд торны нүд бүр нь дотроо байгаа объектын төв цэгийг агуулж байгаа эсэхэд хариуцаж, байршлын координат болон ангиллын магадлалыг таамагладаг. 
\end{itemize}

YOLO алгоритм нь анх танилцуулагдсанаасаа хойш тасралтгүй хөгжиж v1-ээс v8, v9 хүртэл олон хувилбарууд гарсан байдаг. Бидний энэхүү ажилдаа сонгож ашиглаж буй \textbf{YOLOv8} нь 2023 онд Ultralytics багаас гаргасан бөгөөд (өмнөх YOLOv5-ийг зохиосон баг) архитектурын хувьд гүнзгий өөрчлөлтүүд хийгдсэн төгс хувилбаруудын нэг юм \cite{yolov8_ultralytics}. YOLOv8-ийн хамгийн том дэвшил нь anchor-box-т суурилсан аргачлалаас татгалзаж, anchor-free илрүүлэлтийн алгоритмыг (CornerNet, FCOS зэрэгтэй төстэй) ашиглах болсонд оршино. Үүний үр дүнд бараануудын хэлбэр хэмжээнээс хамаарсан гипер-параметрүүдийг (anchor boxes) гараар тохируулах шаардлагагүй болж, загвар сургах явц илүү хялбар болж, бодит амьдрал дээрх том, жижиг янз бүрийн хэмжээтэй объектуудыг таних нарийвчлал өссөн юм. Энэ нь олон төрлийн савлагаа бүхий ундаа, шүүс, шоколад зэрэг жижиглэн худалдааны барааг хагас халхлагдсан байдлаас илрүүлэхэд туйлын тохиромжтой байна.

\section{Жижиглэн худалдааны салбар дахь ухаалаг мониторингийн системүүд}
Жижиглэн худалдаанд зориулсан хиймэл оюун ухаант хяналтын системүүд гадны өндөр хөгжилтэй орнуудад амжилттай нэвтэрч эхлээд байна. Үүний хамгийн тод жишээнүүдийн нэг бол **Amazon Go** ухаалаг дэлгүүр юм \cite{amazon_go}. Amazon Go нь "Just Walk Out" технологиороо дэлгүүрийн таазанд байрлуулсан хэдэн зуун өндөр нягтралтай камерууд болон лангуун дээрх жин хэмжигч мэдрэгчийн тусламжтайгаар худалдан авагчийн авсан бараа бүрийг шууд тооцдог систем юм. Харин уг систем нь бүхэл бүтэн дэд бүтцийг шинээр барих шаардлагатай бөгөөд суурилуулалтын зардал асар өндөр тул зөвхөн тусгайлан баригдсан өндөр зэрэглэлийн дэлгүүрүүдэд л хэрэгждэг хязгаарлалттай.

Дэлгүүрийн менежментэд зориулагдсан өөр нэг алдартай программ хангамжийн шийдэл бол **Trax Retail** юм \cite{trax_retail}. Trax нь дэлгүүрийн ажилчид эсвэл бөөний нийлүүлэлтийн төлөөлөгчид ухаалаг утас болон таблетаараа лангууны зураг дарах бөгөөд тэдгээр зургуудыг клоуд дээр өөрсдийн хаалттай объект таних системээр дамжуулан боловсруулж, хувь хэмжээ, байширлын аудитын тайланг бэлтгэдэг. Энэхүү систем нь манай судалгааны ажилтай үзэл баримтлалын хувьд маш төстэй боловч энэ нь хаалттай эх бүхий үнэтэй арилжааны бүтээгдэхүүн бөгөөд жижиг дунд үйлчилгээний байгууллагууд болон орон нутгийн супермаркетууд ашиглахад үнэтэй тусдаг.

Мөн түүнчлэн, камертай роботууд ашиглан дэлгүүрийн эгнээ дундуур явж тооллого хийдэг **Bossa Nova Robotics** брэндийн шийдлүүдийг Америкийн Walmart сүлжээ дэлгүүрт хэсэгчлэн нэвтрүүлж туршсан түүхтэй. Гэсэн хэдий ч робот өөрөө нарийн гарцтай дэлгүүрийн орчинд үйлчлүүлэгчдэд хүндрэл учруулах, мөн засвар үйлчилгээний зардал өндөр гэсэн шалтгаанаар өргөнөөр нэвтэрч чадаагүй байна. 

\section{Ижил төстэй системүүдийн харьцуулсан шинжилгээ, бидний бүтээлийн байр суурь}
Дээр дурдсан олон төрлийн шийдлүүдтэй харьцуулан үзвэл бидний энэхүү дипломын ажилд хөгжүүлж буй систем нь дараах өвөрмөц онцлог болон давуу талуудтай:

\begin{enumerate}
    \item \textbf{Интеграцийн хялбар байдал болон өртөг хямд:} Системийг ажиллуулахын тулд дэлгүүрт нэмэлт мэдрэгч, үнэтэй камер суурилуулах огт шаардлагагүй бөгөөд борлуулагчдын халаасанд байдаг ухаалаг утсыг гол мэдрэгч (sensor) болгон ашиглана.
    \item \textbf{Нээлттэй эхэд суурилсан дэвшилтэт технологи:} Ultralytics YOLOv8 зэрэг нээлттэй эхийн шилдэг шийдлүүдийг ашиглаж, өөрийн орны зах зээлийн бараа бүтээгдэхүүн (Монгол шошготой, орон нутгийн үйлдвэрийн бараа) дээр тусгайлан загварыг сургасан. Гадны бэлэн системүүд нь Монголын дотоодын барааг таньж чадахгүй байх эрсдэлтэй байдаг бол бидний шийдэл энэ талын датасетийг бэлтгэдгээрээ давуу талтай.
    \item \textbf{Төгсгөл хоорондын (End-to-End) автомат аудит:} Зөвхөн объектыг илрүүлээд гаргаад зогсохгүй, тухайн дэлгүүрийн өгөгдлийн санд бүртгэлтэй байх ёстой стандарт хэмжээ, байширлын зураглал (Planogram)-ийн логиктой харьцуулалт хийж, дүгнэлт статусыг (Аудитын статус) автоматаар зөвлөн харуулдгаараа илүү ухаалаг бөгөөд цогц систем болж чадаж байгаа юм. Энэ нь менежерүүд гараар дүгнэлт хийх шаардлагыг устгаж байгаа хэрэг юм.
\end{enumerate}

Энэхүү судалгаа нь орчин үеийн компьютерын харааны ололт амжилтыг зөвхөн онолын хувьд баталж зогсохүй, бизнесийн бодит хэрэгцээнд хэрхэн хурдан, хялбар, үр ашигтайгаар хөрвүүлэн хэрэглэж болох бодит жишээ болно гэж дүгнэж байна.
